\section{What replaces the guarantee}
\label{sec:certificate}

A method that may fail, and that may fail by stabilising on a set that is simply
not optimal, is usable when the answer is checked. Here
Proposition~\ref{prop:sufficient} carries the argument: a candidate $(x,\lambda)$
satisfying
\begin{equation}
  \norm{G x - a - C\lambda}_\infty \le \varepsilon,
  \quad |c_i\T x - b_i| \le \varepsilon \ (i \le m_{\mathrm{eq}}),
  \quad c_i\T x - b_i \ge -\varepsilon, \ \ \lambda_i \ge -\varepsilon,
  \ \ |\lambda_i (c_i\T x - b_i)| \le \varepsilon
  \label{eq:certificate}
\end{equation}
on the inequalities, with $\varepsilon$ scaled by the data, \emph{is} the answer
to that tolerance. Stationarity is checked against $G$ directly and not inferred
from the construction, since the construction is what an ill-conditioned working
set corrupts. A candidate that fails is discarded and an exact
method runs instead, so guessing can never degrade the answer: it produces the
minimiser or nothing.

This is the trade the rest of the note is about. One gives up finite termination,
which on general constraints was not available anyway, and buys back correctness
by verification. The asymmetry that makes the trade worth taking is that
verifying~\eqref{eq:certificate} costs $O(nm)$ for a dense $C$, one pass over the
constraints, against the $O(n^3)$ of the fallback solve.

Two implementation points follow, and they are the ones a reader who has to build
such a solver will want.

\begin{remark}[The rank test is the factorisation]
There is no need for a separate rank decision. Attempting the Cholesky
factorisation of $C_\Active\T G^{-1}C_\Active$ \emph{is} the test: by
Proposition~\ref{prop:guessrank} the matrix is positive definite exactly when the
guess has full column rank, so a dependent guess fails the factorisation rather
than returning something plausible. What an implementation must not do is solve
the system by a method that tolerates rank deficiency and returns a least-squares
answer, because that converts a detected failure into an undetected one.
\end{remark}

\begin{remark}[Two tolerances with very different jobs]
\label{rem:tolerances}
The tolerance in~\eqref{eq:repair} decides which set is tried next. A bad choice
there costs a repair or a fallback and can never cost correctness, so it is not
delicate. The tolerance in~\eqref{eq:certificate} is the only thing standing
between a non-optimal point and the caller, and is therefore the one to reason
about. It is nonetheless not delicate either, because the quantity it thresholds
is bimodal: a candidate constructed from a well-conditioned working set satisfies
\eqref{eq:certificate} to a few multiples of the machine epsilon, and one built on
a set that is not has no reason to come close. Section~\ref{sec:measurement}
measures the accepted population, whose worst residual is $\qpPdasAcceptWorst$.
\end{remark}

The least-index rule keeps a role under this reading, but a demoted one. It is a
way of making progress where block exchange oscillates, and not a termination
proof. Neither it nor the rank test substitutes for the guarantee. The certificate
does.
